\chapter{Conclusie}%
\label{ch:conclusie}

Deze bachelorproef onderzocht hoe een mobiele applicatie op maat van studenten met ADHD kan bijdragen
aan beter taakbeheer en tijdsplanning, met minder stress en uitstelgedrag. De aanleiding voor dit onderzoek
was het feit dat bestaande planningsapps vaak te complex zijn en onvoldoende rekening houden met de
cognitieve beperkingen die gepaard gaan met ADHD, zoals problemen met executieve functies, tijdsperceptie
en emotieregulatie.

Op basis van een literatuurstudie werd vastgesteld dat vooral extrinsieke cognitieve belasting een belangrijke
factor is bij het falen van klassieke planningstools voor deze doelgroep. Minimalistische interfaces, voorspelbare
navigatie en een focus op kleine stappen kunnen deze belasting reduceren. Daarnaast wijzen studies op het
belang van directe feedback en beloningsmechanismen om motivatie te ondersteunen.

Om deze inzichten te vertalen naar een praktische oplossing werd een proof-of-concept ontwikkeld als een
mobielgericht prototype. De applicatie laat gebruikers toe om taken in te voeren, een specifieke dag te selecteren
voor planning, en voortgang te volgen via een reward systeem met punten en streaks. Daarnaast bevat het prototype
een profielpagina en een instellingenpagina. De implementatie werd uitgevoerd met React, TypeScript en Tailwind CSS,
zonder backend, waardoor de applicatie volledig client-side functioneert.

De centrale onderzoeksvraag kan beantwoord worden door te stellen dat een applicatie op maat van studenten met ADHD
vooral meerwaarde biedt wanneer ze structuur oplegt zonder complexiteit te introduceren. Het prototype toont aan dat een
eenvoudige workflow en directe beloning gecombineerd kunnen worden in een digitale tool die ontworpen is met aandacht
voor cognitieve toegankelijkheid.

% nog toevoegen over gebruikerstesten (moeten nog gebeuren)

Samenvattend kan gesteld worden dat deze bachelorproef een onderbouwde basis biedt voor verdere ontwikkeling van een
ADHD-vriendelijke planningsapplicatie, en dat de combinatie van prikkelarm ontwerp en motiverende feedback een veelbelovende
richting vormt voor digitale ondersteuning binnen het hoger onderwijs.